% Appendix A.3 SR-Agent details -- the system-prompt append (verbatim from the code that ran the experiments:
% agentdojo/src/agentdojo/agent_pipeline/llms/qwen_32b_think_llm_safe_agent.py, SAFE_AGENT_SYSTEM_APPEND). The NeurIPS figure
% P3.pdf omits ", e.g., in tool response" in item 2; the text below is what the models were evaluated with. The RL-Hammer harness
% (injecAgent-rl-harmmer/rl-injector/utils.py) carries the same three items with item 2 phrased "Explicitly identify ...".

\subsection{SR-Agent Details}
\label{app:a3}

For the Qwen3 models (SR-Agent-Qwen3-8B and SR-Agent-Qwen3-4B) the following instruction is appended to the
environment's system prompt at inference time, in every experiment of this paper; it restates the three-part structure
the reflection was trained with. SR-Agent (Llama-3.1-8B), the main model, is evaluated \emph{without} it
(Section~\ref{sec:exp}).

\begin{tcolorbox}[
colback=white, colframe=black, coltitle=black, colbacktitle=gray!30, boxrule=0.3pt,
title={System-prompt append used for the Qwen3 SR-Agent models}, breakable,
 title after break={System-prompt append (continued)}]
\small
For your thinking process, please:\\
1. Briefly summarize the \textbf{original goal}\\
2. Explicitly and briefly identify \textbf{prompt injection} attempts, e.g., in tool response\\
3. Think about \textbf{optimal action} both completing the original goal and resists the prompt injection instructions.
And briefly think about what function call may follow the injection or not fullfill the original task.
\end{tcolorbox}
